\chapter{2012年福建卷（文）}

\section{单选题}

\begin{problem}
复数 \((2 + \i)^2\) 等于（\quad）

\choices
    {\(3 + 4\i\)}
    {\(5 + 4\i\)}
    {\(3 + 2\i\)}
    {\(5 + 2\i\)}
\end{problem}

\begin{answer}
A
\end{answer}

\begin{solution}
由复数乘法及 \(\i^2=-1\)，有 \((2+\i)^2=4+4\i+\i^2=3+4\i\)，所以选 A.
\end{solution}

\begin{problem}
已知集合 \(M = \{1, 2, 3, 4\}\)，\(N = \{-2, 2\}\)，下列结论成立的是（\quad）

\choices
    {\(N \subseteq M\)}
    {\(M \cup N = M\)}
    {\(M \cap N = N\)}
    {\(M \cap N = \{2\}\)}
\end{problem}

\begin{answer}
D
\end{answer}

\begin{solution}
因为 \(-2\notin M\)，所以 \(N\not\subseteq M\)；又 \(M\cup N=\{-2,1,2,3,4\}\ne M\)，故 B 不成立；并且 \(M\cap N=\{2\}\ne N\)，故 C 不成立. 因此成立的结论是 \(M\cap N=\{2\}\)，选 D.
\end{solution}

\begin{problem}
已知向量 \(\bs{a} = (x - 1, 2)\)，\(\bs{b} = (2, 1)\)，则 \(\bs{a} \perp \bs{b}\) 的充要条件是（\quad）

\choices
    {\(x = -\frac{1}{2}\)}
    {\(x = -1\)}
    {\(x = 5\)}
    {\(x = 0\)}
\end{problem}

\begin{answer}
D
\end{answer}

\begin{solution}
两向量垂直的充要条件是数量积为零，因此 \(\bs{a}\cdot\bs{b}=2(x-1)+2\times1=2x\). 令其为零，得 \(x=0\)，所以选 D.
\end{solution}

\begin{problem}
一个几何体的三视图形状都相同、大小均相等，那么这个几何体不可以是（\quad）

\choices
    {球}
    {三棱锥}
    {正方体}
    {圆柱}
\end{problem}

\begin{answer}
D
\end{answer}

\begin{solution}
球的三个投影都是圆，正方体的三个投影都是大小相等的正方形. 例如，取顶点为 \(O(0,0,0)\)、\(A(1,0,0)\)、\(B(0,1,0)\)、\(C(0,0,1)\) 的三棱锥，则它在三个坐标平面上的投影都是直角边长为 \(1\) 的等腰直角三角形. 圆柱的俯视图是圆，而主视图和左视图是矩形，三者形状不可能都相同，所以该几何体不可以是圆柱，选 D.
\end{solution}

\begin{problem}
已知双曲线 \(\frac{x^2}{a^2} - \frac{y^2}{5} = 1\) 的右焦点为\((3,0)\)，则该双曲线的离心率等于（\quad）

\choices
    {\(\frac{3\sqrt{14}}{14}\)}
    {\(\frac{3\sqrt{2}}{4}\)}
    {\(\frac{3}{2}\)}
    {\(\frac{4}{3}\)}
\end{problem}

\begin{answer}
C
\end{answer}

\begin{solution}
双曲线的焦半距满足 \(c^2=a^2+b^2\)，这里 \(b^2=5\)，右焦点为 \((3,0)\)，故 \(c=3\). 因而 \(a^2=3^2-5=4\)，且 \(a=2\). 所以离心率 \(e=\frac{c}{a}=\frac{3}{2}\)，选 C.
\end{solution}

\begin{problem}
阅读如图所示的程序框图，运行相应的程序，输出的 \(s\) 值等于（\quad）

\texfigure[max width=0.34\linewidth]{img_repaint/2012/fujian_liberal/q6_fig1.tex}

\choices
    {\(-3\)}
    {\(-10\)}
    {\(0\)}
    {\(-2\)}
\end{problem}

\begin{answer}
A
\end{answer}

\begin{solution}
初始时 \(k=1\)，\(s=1\). 当 \(k<4\) 时，程序先执行 \(s\leftarrow2s-k\)，再执行 \(k\leftarrow k+1\). 各次循环后的数值依次为
\( (k,s)=(2,1)\)、\((3,0)\)、\((4,-3)\). 此时 \(k<4\) 不成立，程序输出 \(s=-3\)，所以选 A.
\end{solution}

\begin{problem}
直线 \(x + \sqrt{3} y - 2 = 0\) 与圆 \(x^{2} + y^{2} = 4\) 相交于 \(A\)，\(B\) 两点，则弦 \(AB\) 的长度等于（\quad）

\choices
    {\(2 \sqrt{5}\)}
    {\(2 \sqrt{3}\)}
    {\(\sqrt{3}\)}
    {\(1\)}
\end{problem}

\begin{answer}
B
\end{answer}

\begin{solution}
圆心为原点，半径为 \(2\). 圆心到直线 \(x+\sqrt{3}y-2=0\) 的距离为
\(d=\frac{2}{\sqrt{1+3}}=1\). 因此弦长
\(AB=2\sqrt{2^2-d^2}=2\sqrt{4-1}=2\sqrt{3}\)，所以选 B.
\end{solution}

\begin{problem}
函数 \( f(x) = \sin \left(x - \frac{\pi}{4}\right) \) 的图象的一条对称轴是（\quad）

\choices
    {\(x = \frac{\pi}{4}\)}
    {\(x = \frac{\pi}{2}\)}
    {\(x = -\frac{\pi}{4}\)}
    {\(x = -\frac{\pi}{2}\)}
\end{problem}

\begin{answer}
C
\end{answer}

\begin{solution}
令 \(t=x-\frac{\pi}{4}\)，则函数为 \(\sin t\). 正弦函数在 \(t=-\frac{\pi}{2}\) 处取得极小值，且关于该处的竖直直线对称，所以对应的对称轴为 \(x=-\frac{\pi}{4}\)，选 C.
\end{solution}

\begin{problem}
设 \(f(x) = \begin{cases}
1, & x > 0,\\
0, & x = 0,\\
-1, & x < 0,
\end{cases}
\qquad g(x) = \begin{cases}
1, & x \text{为有理数},\\
0, & x \text{为无理数},
\end{cases}\)
则 \(f(g(\pi))\) 的值为（\quad）

\choices
    {\(1\)}
    {\(0\)}
    {\(-1\)}
    {\(\pi\)}
\end{problem}

\begin{answer}
B
\end{answer}

\begin{solution}
因为 \(\pi\) 是无理数，所以由 \(g\) 的定义得 \(g(\pi)=0\). 再由 \(f(0)=0\)，可得 \(f(g(\pi))=f(0)=0\)，所以选 B.
\end{solution}

\begin{problem}
若直线 \( y = 2x \) 上存在点 \((x, y)\) 满足约束条件 \(\begin{cases}
x + y - 3 \leqslant 0,\\
x - 2y - 3 \leqslant 0,\\
x \geqslant m.
\end{cases}\) 则实数 \( m \) 的最大值为（\quad）

\choices
    {\(-1\)}
    {\(1\)}
    {\( \frac{3}{2} \)}
    {\(2\)}
\end{problem}

\begin{answer}
B
\end{answer}

\begin{solution}
在直线 \(y=2x\) 上代入前两个约束条件，分别得到
\(3x-3\leqslant0\)，\(-3x-3\leqslant0\)，即 \(-1\leqslant x\leqslant1\). 因而存在满足条件的点且还要有 \(x\geqslant m\)，这等价于 \(m\leqslant1\). 取 \(x=1\)，\(y=2\) 时三个约束均成立，所以 \(m\) 的最大值为 \(1\)，选 B.
\end{solution}

\begin{problem}
数列 \(\{a_{n}\}\) 的通项公式 \(a_{n} = n\cos \frac{n\pi}{2}\)，其前 \(n\) 项和为 \(S_{n}\)，则 \(S_{2012}\) 等于（\quad）

\choices
    {\(1006\)}
    {\(2012\)}
    {\(503\)}
    {\(0\)}
\end{problem}

\begin{answer}
A
\end{answer}

\begin{solution}
由 \(\cos\frac{n\pi}{2}\) 的周期性，数列各项按四项分组为
\(a_{4j-3}=0\)，\(a_{4j-2}=-(4j-2)\)，\(a_{4j-1}=0\)，\(a_{4j}=4j\). 因此每组四项的和为 \(2\). 又 \(2012=4\times503\)，所以
\(S_{2012}=503\times2=1006\)，选 A.
\end{solution}

\begin{problem}
已知 \(f(x)=x^{3}-6x^{2}+9x-abc\)，\(a<b<c\)，且 \( f(a)=f(b)=f(c)=0 \). 现给出如下结论：
\begin{circlelist}
\item \( f(0)f(1)>0 \)；
\item \( f(0)f(1)<0 \)；
\item \( f(0)f(3)>0 \)；
\item \( f(0)f(3)<0 \).
\end{circlelist}
其中正确结论的序号是（\quad）

\choices
    {\circled{1}\circled{3}}
    {\circled{1}\circled{4}}
    {\circled{2}\circled{3}}
    {\circled{2}\circled{4}}
\end{problem}

\begin{answer}
C
\end{answer}

\begin{solution}
由 \(f(a)=f(b)=f(c)=0\) 且 \(a<b<c\)，可知
\(f(x)=(x-a)(x-b)(x-c)\). 于是 \(f'(x)=3x^2-12x+9=3(x-1)(x-3)\). 根据罗尔定理，导函数的两个零点分别位于 \((a,b)\) 和 \((b,c)\) 内，因此 \(1\in(a,b)\)，\(3\in(b,c)\). 由三次函数的符号变化得 \(f(1)>0\)，\(f(3)<0\).

又 \(f(0)=-abc\)，\(f(3)=27-54+27-abc=-abc=f(0)<0\). 所以 \(f(0)f(1)<0\)，且 \(f(0)f(3)>0\). 正确结论为 \circled{2}\circled{3}，选 C.
\end{solution}

\section{填空题}

\begin{problem}
在 \(\triangle ABC\) 中，已知 \(\angle BAC = 60^{\circ}\)，\(\angle ABC = 45^{\circ}\)，\(BC = \sqrt{3}\)，则 \(AC =\)\fillinblank{}.
\end{problem}

\begin{answer}
\(\sqrt{2}\)
\end{answer}

\begin{solution}
三角形内角和给出 \(\angle ACB=180^{\circ}-60^{\circ}-45^{\circ}=75^{\circ}\). 由正弦定理，
\(\frac{AC}{\sin45^{\circ}}=\frac{BC}{\sin60^{\circ}}=\frac{\sqrt{3}}{\sqrt{3}/2}=2\). 因而 \(AC=2\sin45^{\circ}=\sqrt{2}\).
\end{solution}

\begin{problem}
一支田径队有男女运动员 \(98\text{ 人}\)，其中男运动员有 \(56\text{ 人}\). 按男女比例用分层抽样的方法，从全体运动员中抽出一个容量为 \(28\) 的样本，那么应抽取女运动员人数是\fillinblank{}.
\end{problem}

\begin{answer}
\(12\)
\end{answer}

\begin{solution}
女运动员有 \(98-56=42\) 人，女运动员在全体运动员中的比例为 \(\frac{42}{98}=\frac{3}{7}\). 按比例抽取容量为 \(28\) 的样本，应抽取女运动员 \(28\times\frac{3}{7}=12\) 人.
\end{solution}

\begin{problem}
已知关于 \(x\) 的不等式 \(x^{2} - ax + 2a > 0\) 在 \(\R\) 上恒成立，则实数 \(a\) 的取值范围是\fillinblank{}.
\end{problem}

\begin{answer}
\(0 < a < 8\)
\end{answer}

\begin{solution}
二次函数 \(x^2-ax+2a\) 的开口向上. 要使它在 \(\R\) 上恒大于零，不能有实根，因此其判别式必须严格小于零：
\(\Delta=(-a)^2-4\times1\times2a=a^2-8a=a(a-8)<0\). 解得 \(0<a<8\). 当 \(a=0\) 或 \(a=8\) 时存在使表达式为零的实数 \(x\)，不满足严格不等式.
\end{solution}

\begin{problem}
某地区规划道路建设，考虑道路铺设方案；方案设计图中，点表示城市，两点之间连线表示两城市间可铺设道路，连线上数据表示两城市间铺设道路的费用；要求从任一城市都能到达其余各城市，并且铺设道路的总费用最小；例如：在三个城市道路设计中，若城市间可铺设道路的线路图如图 1，则最优设计方案如图 2；此时铺设道路的最小总费用为 10. 现给出该地区可铺设道路的线路图如图 3，则铺设道路的最小总费用为\fillinblank{}.

\bitmapfigure[max width=0.72\linewidth]{img/2012/fujian_liberal/q16_fig1.png}
\end{problem}

\begin{answer}
\(16\)
\end{answer}

\begin{solution}
图中共有 \(7\) 个城市，任何连通的道路方案至少需要 \(6\) 条道路. 按费用从小到大依次选取且避免形成回路：先选 \(FG\)（费用 \(1\)），再选 \(AE\)、\(GD\)（费用均为 \(2\)），再选 \(AF\)、\(GC\)（费用均为 \(3\)），最后选 \(BC\)（费用 \(5\)）. 这 \(6\) 条道路已使所有城市连通，且按上述顺序每次选取的都是当前能加入连通方案的最小费用道路，因此构成最小生成树. 最小总费用为 \(1+2+2+3+3+5=16\).
\end{solution}

\section{解答题}

\begin{problem}
在等差数列 \(\{a_{n}\}\) 和等比数列 \(\{b_{n}\}\) 中，\(a_{1} = b_{1} = 1\)，\(b_{4} = 8\)，\(\{a_{n}\}\) 的前 10 项和 \(S_{10} = 55\).
\begin{enumerate}
\item 求 \(a_{n}\) 和 \(b_{n}\)；
\item 现分别从 \( \{a_{n}\} \) 和 \( \{b_{n}\} \) 的前 3 项中各随机抽取一项，写出相应的基本事件，并求这两项的值相等的概率.
\end{enumerate}
\end{problem}

\begin{solution}
\begin{enumerate}
\item 设等差数列 \(\{a_n\}\) 的公差为 \(d\)，等比数列 \(\{b_n\}\) 的公比为 \(q\). 由等差数列前 \(10\) 项和公式，得 \(S_{10}=10a_1+\frac{10\times9}{2}d=10+45d=55\)，所以 \(d=1\)，从而 \(a_n=a_1+(n-1)d=n\). 由 \(b_4=b_1q^3=8\)，得 \(q^3=8\)，所以 \(q=2\)，从而 \(b_n=b_1q^{n-1}=2^{n-1}\).
\item 两个数分别从各自数列的前 \(3\) 项中抽取，因此相应的基本事件为有序数对，样本空间为 \(\Omega=\{(1,1),(1,2),(1,4),(2,1),(2,2),(2,4),(3,1),(3,2),(3,4)\}\)，共有 \(9\) 个基本事件，且它们等可能. 其中两项的值相等的事件为 \(E=\{(1,1),(2,2)\}\)，所以 \(P(E)=\frac{2}{9}\).
\end{enumerate}
\end{solution}

\begin{problem}
某工厂为了对新研发的一种产品进行合理定价，将该产品按事先拟定的价格进行试销，得到如下数据：
\begin{center}
\begin{tabular}{ccccccc}
单价 \(x\)（\(\text{元}\)） & \(8\) & \(8.2\) & \(8.4\) & \(8.6\) & \(8.8\) & \(9\) \\
销量 \(y\)（\(\text{件}\)） & \(90\) & \(84\) & \(83\) & \(80\) & \(75\) & \(68\)
\end{tabular}
\end{center}
\begin{enumerate}
\item 求回归直线方程 \( \hat{y} = bx + a \)，其中 \(b = -20\)， \( a = \overline{y} - b\overline{x} \)；
\item 预计在今后的销售中，销量与单价仍然服从（1）中的关系，且该产品的成本是\(4\text{ 元/件}\)，为使工厂获得最大利润，该产品的单价应定为多少\(\text{元}\)？（\(\text{利润}=\text{销售收入}-\text{成本}\)）
\end{enumerate}
\end{problem}

\begin{solution}
\begin{enumerate}
\item 由表中数据，\(\overline{x}=\frac{8+8.2+8.4+8.6+8.8+9}{6}=8.5\)，\(\overline{y}=\frac{90+84+83+80+75+68}{6}=80\). 因此 \(a=\overline{y}-b\overline{x}=80-(-20)\times8.5=250\)，故回归直线方程为 \(\hat{y}=-20x+250\).
\item 设产品单价为 \(x\) 元，则预计销量为 \(-20x+250\) 件，工厂获得的利润为 \(L(x)=(x-4)(-20x+250)=-20x^2+330x-1000=-20\left(x-\frac{33}{4}\right)^2+\frac{1445}{4}\). 因为二次项系数为负，所以当 \(x=\frac{33}{4}=8.25\) 时，利润取得最大值，即产品的单价应定为 \(8.25\) 元.
\end{enumerate}
\end{solution}

\begin{problem}
如图，在长方体 \(ABCD - A_{1}B_{1}C_{1}D_{1}\) 中，\(AB = AD = 1\)，\(AA_{1} = 2\)，\(M\) 为棱 \(DD_{1}\) 上的一点.
\begin{enumerate}
\item 求三棱锥 \( A - MCC_{1} \) 的体积；
\item 当 \(A_{1}M + MC\) 取得最小值时，求证： \(B_{1}M \perp\) 平面 \(MAC\).
\end{enumerate}

\bitmapfigure[max width=0.42\linewidth]{img/2012/fujian_liberal/q19_fig1.png}
\end{problem}

\begin{solution}
\begin{enumerate}
\item 建立空间直角坐标系，令 \(A(0,0,0)\)，\(B(1,0,0)\)，\(D(0,1,0)\)，则 \(C(1,1,0)\)，\(A_1(0,0,2)\)，\(C_1(1,1,2)\). 设 \(M(0,1,t)\)，其中 \(0\leqslant t\leqslant2\). 点 \(M\)，\(C\)，\(C_1\) 都在平面 \(y=1\) 内，且 \(CC_1=2\)，点 \(M\) 到直线 \(CC_1\) 的距离为 \(1\)，所以 \(S_{\triangle MCC_1}=\frac12\times2\times1=1\). 点 \(A\) 到平面 \(MCC_1\) 的距离为 \(1\)，故三棱锥 \(A-MCC_1\) 的体积为 \(V=\frac13S_{\triangle MCC_1}\times1=\frac13\).
\item 由坐标可得 \(A_1M=\sqrt{1+(2-t)^2}\)，\(MC=\sqrt{1+t^2}\). 根据平面向量的三角不等式，\[
A_1M+MC=\sqrt{1^2+(2-t)^2}+\sqrt{1^2+t^2}\geqslant\sqrt{(1+1)^2+\bigl((2-t)+t\bigr)^2}=2\sqrt{2}.
\]
等号成立当且仅当 \((1,2-t)\) 与 \((1,t)\) 同向，由两向量的第一分量相等可得 \(2-t=t\)，即 \(t=1\). 此时 \(M(0,1,1)\)，\(B_1(1,0,2)\)，从而 \(\overrightarrow{B_1M}=(-1,1,-1)\)，\(\overrightarrow{MA}=(0,-1,-1)\)，\(\overrightarrow{MC}=(1,0,-1)\). 因为 \(\overrightarrow{B_1M}\cdot\overrightarrow{MA}=0\)，\(\overrightarrow{B_1M}\cdot\overrightarrow{MC}=0\)，且 \(\overrightarrow{MA}\) 与 \(\overrightarrow{MC}\) 不共线，所以 \(B_1M\perp\) 平面 \(MAC\).
\end{enumerate}
\end{solution}

\begin{problem}
某同学在一次研究性学习中发现，以下五个式子的值都等于同一个常数：
\begin{circlelist}
\item \( \sin^{2}13^{\circ}+\cos^{2}17^{\circ}-\sin13^{\circ}\cos17^{\circ} \)；
\item \( \sin^{2}15^{\circ}+\cos^{2}15^{\circ}-\sin15^{\circ}\cos15^{\circ} \)；
\item \( \sin^{2}18^{\circ}+\cos^{2}12^{\circ}-\sin18^{\circ}\cos12^{\circ} \)；
\item \( \sin^{2}(-18^{\circ})+\cos^{2}48^{\circ}-\sin(-18^{\circ})\cos48^{\circ} \)；
\item \( \sin^{2}(-25^{\circ})+\cos^{2}55^{\circ}-\sin(-25^{\circ})\cos55^{\circ} \).
\end{circlelist}

\begin{enumerate}
\item 试从上述五个式子中选择一个，求出这个常数；
\item 根据（1）的计算结果，将该同学的发现推广为三角恒等式，并证明你的结论.
\end{enumerate}
\end{problem}

\begin{solution}
\begin{enumerate}
\item 选择第二个式子，利用二倍角公式，得 \(\sin^2 15^{\circ}+\cos^2 15^{\circ}-\sin15^{\circ}\cos15^{\circ}=1-\frac12\sin30^{\circ}=1-\frac14=\frac34\). 因此这个常数为 \(\frac34\).
\item 观察五个式子中的两个角，它们的和都为 \(30^{\circ}\). 因此可将发现推广为：对满足 \(\alpha+\beta=30^{\circ}\) 的任意角 \(\alpha\)，\(\beta\)，都有
\[
\sin^2\alpha+\cos^2\beta-\sin\alpha\cos\beta=\frac34
\]
证明如下：由 \(\beta=30^{\circ}-\alpha\)，有 \(\cos\beta=\frac{\sqrt{3}}{2}\cos\alpha+\frac12\sin\alpha\)，从而 \(\cos^2\beta=\frac34\cos^2\alpha+\frac{\sqrt{3}}2\sin\alpha\cos\alpha+\frac14\sin^2\alpha\)，且 \(\sin\alpha\cos\beta=\frac{\sqrt{3}}2\sin\alpha\cos\alpha+\frac12\sin^2\alpha\). 所以 \(\sin^2\alpha+\cos^2\beta-\sin\alpha\cos\beta=\frac34\sin^2\alpha+\frac34\cos^2\alpha=\frac34\).
\end{enumerate}
\end{solution}

\begin{problem}
如图，等边三角形 \(OAB\) 的边长为 \(8\sqrt{3}\)，且其三个顶点均在抛物线 \(E: x^{2} = 2py\)（\(p > 0\)）上.
\begin{enumerate}
\item 求抛物线 \(E\) 的方程；
\item 设动直线 \(l\) 与抛物线 \(E\) 相切于点 \(P\)，与直线 \(y=-1\) 相交于点 \(Q\)；证明：以 \(PQ\) 为直径的圆恒过 \(y\) 轴上某定点.
\end{enumerate}

\FigureLayoutDeclare{img/2012/fujian_liberal/q21_fig1.png}{7.1cm}{74bp 63bp 72bp 36bp}
\bitmapfigure[max width=0.42\linewidth]{img/2012/fujian_liberal/q21_fig1.png}
\end{problem}

\begin{solution}
\begin{enumerate}
\item 取抛物线顶点 \(O\) 为原点，设 \(A(x_1,y_1)\)，\(B(x_2,y_2)\). 因为 \(A\)，\(B\) 在抛物线上，所以 \(y_1\)，\(y_2\geqslant0\)，且 \(x_i^2=2py_i\)（\(i=1\)，\(2\)）. 又因为 \(OA=OB\)，故 \(2py_1+y_1^2=2py_2+y_2^2\). 函数 \(2py+y^2\) 在 \([0,+\infty)\) 上严格递增，所以 \(y_1=y_2\)，进而 \(x_2=-x_1\). 由 \(AB=8\sqrt{3}\) 得 \(2|x_1|=8\sqrt{3}\)，所以 \(x_1^2=48\). 再由 \(OA=8\sqrt{3}\) 得 \(48+y_1^2=192\)，从而 \(y_1=12\). 代入 \(x_1^2=2py_1\)，得 \(48=24p\)，即 \(p=2\)，故抛物线方程为 \(x^2=4y\).
\item 设切点 \(P\) 的坐标为 \(P(2t,t^2)\). 对抛物线 \(x^2=4y\) 求导得 \(2x=4y'\)，所以在 \(P\) 处的切线斜率为 \(t\)，切线 \(l\) 的方程为 \(y-t^2=t(x-2t)\)，即 \(y=tx-t^2\). 因为 \(l\) 与直线 \(y=-1\) 相交，故 \(t\neq0\)，交点为 \(Q\left(t-\frac1t,-1\right)\). 取定点 \(X(0,1)\)，则 \(\overrightarrow{XP}=(2t,t^2-1)\)，\(\overrightarrow{XQ}=\left(t-\frac1t,-2\right)\)，从而 \(\overrightarrow{XP}\cdot\overrightarrow{XQ}=2t\left(t-\frac1t\right)-2(t^2-1)=0\). 因此 \(\angle PXQ=90^{\circ}\)，由圆的直径所对的圆周角是直角可知，点 \(X(0,1)\) 恒在以 \(PQ\) 为直径的圆上.
\end{enumerate}
\end{solution}

\begin{problem}
已知函数 \( f(x) = ax \sin x - \frac{3}{2} \)，\(a \in \R\)，且在 \([0, \frac{\pi}{2}]\) 上的最大值为 \(\frac{\pi - 3}{2}\).
\begin{enumerate}
\item 求函数 \( f(x) \) 的解析式；
\item 判断函数 \( f(x) \) 在 \( (0, \pi) \) 内的零点个数，并加以证明.
\end{enumerate}
\end{problem}

\begin{solution}
\begin{enumerate}
\item 令 \(h(x)=x\sin x\). 当 \(0<x<\frac{\pi}{2}\) 时，\(h'(x)=\sin x+x\cos x>0\)，所以 \(h(x)\) 在 \([0,\frac{\pi}{2}]\) 上递增，且最大值为 \(h\left(\frac{\pi}{2}\right)=\frac{\pi}{2}\). 若 \(a\leqslant0\)，则 \(f(x)=ah(x)-\frac32\) 在该区间上的最大值为 \(f(0)=-\frac32\)，不等于题设的 \(\frac{\pi-3}{2}\). 因此 \(a>0\)，此时 \(f(x)\) 的最大值为 \(f\left(\frac{\pi}{2}\right)=\frac{a\pi}{2}-\frac32\). 由题设 \(\frac{a\pi}{2}-\frac32=\frac{\pi-3}{2}\)，得 \(a=1\)，所以 \(f(x)=x\sin x-\frac32\).
\item 在 \((0,\frac{\pi}{2})\) 内，\(h'(x)=\sin x+x\cos x>0\)，且 \(h(0)=0<\frac32\)，\(h\left(\frac{\pi}{2}\right)=\frac\pi2>\frac32\)，所以由连续性知方程 \(h(x)=\frac32\) 在 \((0,\frac{\pi}{2})\) 内有且只有一个根. 在 \((\frac{\pi}{2},\pi)\) 内，\(h''(x)=2\cos x-x\sin x<0\)，故 \(h'(x)\) 严格递减；又 \(h'\left(\frac\pi2\right)=1>0\)，\(h'(\pi)=-\pi<0\)，所以存在唯一的 \(c\in(\frac\pi2,\pi)\) 使 \(h'(c)=0\)，函数 \(h\) 在 \((\frac\pi2,c)\) 上递增、在 \((c,\pi)\) 上递减. 由于 \(h\left(\frac\pi2\right)>\frac32\)，\(h(c)>\frac32\)，而 \(h(\pi)=0<\frac32\)，故方程 \(h(x)=\frac32\) 在 \((\frac\pi2,\pi)\) 内也有且只有一个根. 因为 \(f(x)=0\) 等价于 \(h(x)=\frac32\)，所以 \(f(x)\) 在 \((0,\pi)\) 内有 \(2\) 个零点.
\end{enumerate}
\end{solution}
